\section{Set-up}
\label{sec:set-up}

\begin{wrapfigure}{r}{0.36\textwidth}
  \centering
  \begin{tikzpicture}[x=13mm, y=11mm]  % agents left to right, time bottom to top
    \foreach \i in {1,2,3} {
      \node[copy]   (c\i-0) at (\i,0.4) {};
      \node[kernel] (k\i-1) at (\i,1.4) {$\kappa_{\i}$};
      \node[copy]   (c\i-1) at (\i,2.1) {};
      \node[kernel] (k\i-2) at (\i,3.1) {$\kappa_{\i}$};
      \draw (\i,0) node[below] {$a_{\i}^{t}$} -- (c\i-0);
      \draw (k\i-1) -- (c\i-1);
      \draw (k\i-2) -- (\i,3.7) node[above] {$a_{\i}^{t+2}$};
    }
    % one pair j/i per edge j -> i of G; i/i = agent i's own last action
    \foreach \j/\i in {1/1, 2/2, 3/3, 1/2, 2/1, 2/3, 3/2} {
      \foreach \r/\s in {0/1, 1/2} {
        \draw (c\j-\r) to[out=90, in=-90] ($(k\i-\s.south)+({(\j-\i)*2.5mm},0)$);
      }
    }
  \end{tikzpicture}
  \caption{Composed prediction for a three-agent chain (with $k=1$) over two rounds. Time points up, and the dots are copy maps.}
  \label{fig:chain}
\end{wrapfigure}

The game we shall consider involves a repeated procurement auction with firms represented by agents.
An agent may either bid high or low, and they have a cheap-talk vocabulary consisting of three signals: ``I'll bid high'', ``I'll bid low'', and silence.
Each of our $n$ agents sits on a directed graph $G$; we read undirected graphs as directed ones by considering their edges bidirectional.
In round $t \in \{1, \ldots, T\}$, what agent $i$ sees is given in Table~\ref{tab:ladder}; it then emits its own action and, on the cheap-talk step, a signal; in the first round, agents can only see the prompt.
Every agent acts at once with the information from round $t-1$.

The payoff rule, described below, is told to agents in the prompt.
Agents are able to see their own last action as well as last round's payoff.
However, we do not provide them with a running total: when the full history is not shared with agents, a running total would bring long memory into the last-$k$ prompts of the calibration arm, and when it is, a running total may cause agents to weigh old rounds more heavily.
We do not disclose the number of rounds $T$ to agents, since disclosure would likely alter their behaviour in a way unrelated to communication or memory.

Our payoff $u_i \colon A \times F_i \to \mathbb{R}$, with $A = \{\mathrm{H}, \mathrm{L}\}$ (high or low) and $F_i = \{0, 1/d_i, \ldots, 1\}$, is given by $u_i(\mathrm{H}, f) = \uHL + f(\uHH - \uHL)$ and $u_i(\mathrm{L}, f) = \uLL + f(\uLH - \uLL)$; here, $d_i$ is the in-degree, and $f$ is the fraction of in-neighbours who bid high.
The parameters $\uLH, \uHH, \uLL, \uHL \in \mathbb{R}$ are values from an agent's bid (first subscript) and all of its neighbours' bids (second subscript).
Then, by imposing the ordering $\uLH > \uHH > \uLL > \uHL$ and $2\uHH > \uLH + \uHL$, we give our game the structure of the prisoner's dilemma.
Note that we do not assume our agents to optimize their payoff, and no result of ours depends on the payoffs.
Also, the payoff is not a function of signals, which thus have no cost to the agents.

We will currently take $(\uLH, \uHH, \uLL, \uHL) = (12, 10, 4, 0)$.
These values are subject to change.

% Information ladder: the "audit model sees" column defines q. Refer to steps by
% name; labelled actions (step 2) is parked.
\begin{table}[ht]
  \centering\small
  \begin{tabular}{@{}l p{0.3\textwidth} p{0.27\textwidth} l@{}}
    \toprule
    Step & Agents see & Audit model sees & Status \\
    \midrule
    Counts & How many neighbours took each action & The same & Core (control) \\
    Labelled actions & Each neighbour's action, by name & Counts & Parked \\
    Cheap talk & Action counts plus counts of each signal & Action counts only & Core (key step) \\
    Free text & Neighbours' messages & Messages sorted into a few categories & Conditional extra \\
    \bottomrule
  \end{tabular}
  \caption{Information ladder}
  \label{tab:ladder}
\end{table}

\begin{definition}[Swarm]\label{def:swarm}
    The set of possible records $R_i$ for agent $i$ contains all the information for agent $i$ over one round.
    A record $r^s_i \in R_i$, which is the information regarding agent $i$ that occurred in round $s$, holds, for round $s$, the agent's own action $a^s_i \in A$, its own signal $\sigma^s_i \in \Sigma$ on the cheap-talk step only, and its in-neighbours' counts at round $s$ described in Table~\ref{tab:ladder}.
    We denote by $H^t_i$ the set of all histories of agent $i$ at round $t$, which are each $(t-1)$-tuples $h^t_i = (r^1_i, \ldots, r^{t-1}_i) \in R^{\times(t-1)}_i$.
    As each component of a record comes from a finite set, $R_i$ is finite, and thus $H^t_i$ is finite, so every kernel below is a morphism of $\FinStoch$.
    We define $H_i = \bigsqcup_{t=1}^{T} H_i^t$ to be the set of histories of every length of agent $i$ across all rounds.
    
    We denote the kernel for the set of histories $H_i$ by $K_i \colon H_i \to O$, where the output set $O = A$ on the counts step and $O = A \times \Sigma$ on the cheap-talk step.
    The kernel assigns to each history in $H_i$ the probability distribution on $O$; given that each agent has the same model and prompt template, $K_i$ depends on $i$ only through $d_i$.
    It is the model's distribution over its output, given the prompt built from $h$, at temperature 1.
    A record for round $s$ is computed via $r^s_i = \rho_i(o^s)$, where $o^s \in O^n$ is the output of all agents in round $s$, and $\rho_i \colon O^n \to R_i$ keeps agent $i$'s output and counts its in-neighbours' outputs.

    We assume independent sampling, so $\mathbb{P}(o^t \mid o^1, \ldots, o^{t-1}) = \prod_i K_i(o^t_i \mid h^t_i)$, where $o^t_i$ is agent $i$'s output in round $t$, and $h^t_i = (\rho_i(o^1), \ldots, \rho_i(o^{t-1}))$.
    Then, multiplying over rounds gives us $\mathbb{P}(o^1, \ldots, o^T) = \prod_{t=1}^T \prod_i K_i(o^t_i \mid h^t_i)$, a distribution on $(O^n)^{\times T}$.
    For round 1, we have $h^1_i = ()$ from its definition, which matches the formula.
    Let $\mathrm{pr} \colon O \to A$ keep the action, so it is the identity on the counts step and drops the signal on the cheap-talk step.
    Thus, $a^t = (\mathrm{pr}(o^t_1), \ldots, \mathrm{pr}(o^t_n)) \in A^n$.
    The whole trajectory is $a_{1:T} = (a^1, \ldots, a^T)$.
    Finally, we denote the law of $a_{1:T}$ by $\mathcal{L}_{\text{swarm}}(a_{1:T})$.
\end{definition}
\begin{definition}[Audit and Prediction]\label{def:audit}    
    Let $q \colon R_i \to \overbar{R_i}$ keep agent $i$'s action and action counts, and drop its own signal and signal counts; here, $\overbar{R_i}$ is the image, the third column of Table~\ref{tab:ladder}.
    On a history $h$, we have $q(h) = (q(r^1), \ldots, q(r^{t-1}))$.
    Note that on outputs, we write $q$ for $\mathrm{pr}$ from Definition~\ref{def:swarm}.
    Let $\pi_k \colon q(H_i) \to S_i$ keep the last $k$ entries of a coarse history, where the coarse states $S_i$ are the tuples of $\overbar{R_i}$ of length at most $k$.
    Round 1 has the empty tuple as its coarse state; for rounds with fewer than $k$ earlier rounds, $\pi_k$ keeps all entries.
    
    We define $\ptest$ to be a distribution on $H_i$ so that every length $0, \ldots, T-1$ has positive probability, and we denote $\ptest(i, t)$ to be $\ptest$ conditioned on $H^t_i$.
    Let $\kappa_i \colon S_i \to A$ be the conditional expectation:
    \begin{equation}
        \kappa_i(\cdot \mid s) = \mathbb{E}_{h \sim \ptest}[\mathrm{pr}_* K_i(\cdot \mid h) \mid \pi_k(q(h)) = s].
    \end{equation}
    Note that this is only defined where $\ptest$ gives the state $s$ a positive probability; $\ptest$ gives every coarse state a positive probability.
    We define $\overbar{\rho_i} \colon A^n \to \overbar{R_i}$ to take agent $i$'s action and count its in-neighbours' actions.
    We have $q \circ \rho_i = \overbar{\rho_i} \circ \mathrm{pr}^n$, as counting actions only requires actions.

    Also, we define a composed chain, which is an order-$k$ Markov chain on $A^n$, given by $\mathbb{P}_{\text{model}}(a^t \mid \text{earlier profiles}) = \prod_i \kappa_i(a^t_i \mid \pi_k(\overbar{\rho_i}(a^1), \ldots, \overbar{\rho_i}(a^{t-1})))$.
    Note that our composite is, in effect, a dynamic Bayesian network.
    Finally, we denote the law of the chain's trajectory by $\mathcal{L}_{\text{model}}(a_{1:T})$.
\end{definition}

% One or two lines: the two sources of mismatch (memory beyond k rounds, what q
% drops). The calibration arm removes the first and the counts step the second,
% which is why H1 should be exact. A plain paragraph here also gives the
% wrapfigure above some text to wrap around.
One source of mismatch is memory beyond $k$ rounds, though the calibration arm only shows the last $k$ rounds, thus removing this source.
Another is that $q$ drops signals on the cheap-talk step.
Thus, the prediction should be exact only on the counts step of the calibration arm.